\documentclass[a4paper,11pt]{article}
\usepackage{exam-style-341}
\examinehead{341 农业综合知识三 · 程序设计}{模拟试卷（一）}
\begin{document}
\examcover{341 农业综合知识三}{程序设计部分 · 模拟试卷（一）}{程序设计（C 语言）（50 分）}{50 分}{60 分钟}{本科目只考程序设计，与数据库、网络无关。}

\parttitle{程序设计（C 语言）}

\qtypename{一、单项选择题}{10 题，每题 2 分，共 20 分}
\begin{qgroup}
\q 下列选项中，合法的 C 语言标识符是（\quad）。
\choicesline{A. \code{\_num2}}{B. \code{2num}}{C. \code{-num}}{D. \code{double}}

\q 下列运算符中优先级最低的是（\quad）。
\choicesline{A. \code{!}}{B. \code{\&\&}}{C. \code{=}}{D. \code{!=}}

\q 执行下列程序段后，变量 \code{k} 的值是（\quad）。
\begin{lstlisting}[language={[custom]C}]
int k = 3;
switch (k) {
    case 3: k++;
    case 4: k++; break;
    case 5: k++;
    default: k++;
}
\end{lstlisting}
\choicesline{A. 4}{B. 5}{C. 6}{D. 7}

\q 下列程序段的输出结果是（\quad）。
\begin{lstlisting}[language={[custom]C}]
int i, s = 0;
for (i = 1; i <= 5; i++) {
    if (i % 2 == 0) continue;
    s += i;
}
printf("%d", s);
\end{lstlisting}
\choicesline{A. 6}{B. 9}{C. 15}{D. 5}

\q 若有定义 \code{char s[] = "China";}，则 \code{sizeof(s)} 与 \code{strlen(s)} 的值分别是（\quad）。
\choicesline{A. 6 和 5}{B. 5 和 5}{C. 6 和 6}{D. 5 和 6}

\q 下列程序的输出结果是（\quad）。
\begin{lstlisting}[language={[custom]C}]
#include <stdio.h>
void f(int x) { x = x + 10; }
int main(void) {
    int a = 5;
    f(a);
    printf("%d", a);
    return 0;
}
\end{lstlisting}
\choicesline{A. 5}{B. 10}{C. 15}{D. 0}

\q 若有定义 \code{int a[5] = \{1, 3, 5, 7, 9\}, *p = a;}，则表达式 \code{*(p + 2) + p[4]} 的值是（\quad）。
\choicesline{A. 12}{B. 14}{C. 16}{D. 9}

\q 关于结构体和共用体，下列说法中正确的是（\quad）。
\choices{A. 共用体变量的各成员共用同一段存储单元，该变量所占字节数等于最长成员所占的字节数}{B. 结构体变量的各成员也共用同一段存储单元}{C. 共用体变量可以在定义时对全部成员同时赋初值}{D. 结构体变量不能作为函数的参数}

\q 若 \code{fp} 已定义为 \code{FILE *} 类型，要以“只读”方式打开当前目录下已存在的文本文件 \code{data.txt}，正确的语句是（\quad）。
\choices{A. \code{fp = fopen("data.txt", "r");}}{B. \code{fp = fopen("data.txt", "w");}}{C. \code{fp = fopen("data.txt", "a");}}{D. \code{fp = fopen("data.txt", "rb+");}}

\q 若有定义 \code{char a[20] = "Hello", b[] = "World";}，执行下列语句后的输出结果是（\quad）。
\begin{lstlisting}[language={[custom]C}]
strcat(a, b);
printf("%d", strlen(a));
\end{lstlisting}
\choicesline{A. 5}{B. 10}{C. 11}{D. 20}
\end{qgroup}

\qtypename{二、填空题}{5 题，每题 2 分，共 10 分}
\begin{qgroup}
\q 若有定义 \code{int a = 7, b = 3; double c = 2.0;}，则表达式 \code{a / b} 的值是 \fillblank{}，表达式 \code{a / c} 的值是 \fillblank{}。

\q 用 C 语言表达式描述“\code{x} 是能被 7 整除的偶数”：\fillblank{}。

\q 若有定义 \code{char str[] = "exam";}，则 \code{strlen(str)} 的值是 \fillblank{}，\code{str[1]} 的值是 \fillblank{}。

\q 已知函数定义为 \code{void swap(int *p, int *q) \{ int t = *p; *p = *q; *q = t; \}}，主函数中已有 \code{int x = 1, y = 2;}，执行 \code{swap(\&x, \&y);} 之后，\code{x} 的值是 \fillblank{}，\code{y} 的值是 \fillblank{}。

\q 若已定义 \code{FILE *fp;}，则以只读方式打开当前目录下文本文件 \code{in.txt} 的语句是 \code{fp = fopen(\fillblanks[2.4cm], "r");}，读写结束后应调用 \fillblank{} 函数关闭该文件。
\end{qgroup}

\qtypename{三、程序阅读题}{2 题，每题 5 分，共 10 分}
\begin{qgroup}
\q （5 分）写出下列程序的运行结果。
\begin{lstlisting}[language={[custom]C}]
#include <stdio.h>
int fun(int *p, int n) {
    int i, s = 0;
    for (i = 0; i < n; i++)
        if (p[i] % 2 != 0) s += p[i];
    return s;
}
int main(void) {
    int a[8] = {3, 1, 4, 1, 5, 9, 2, 6};
    int i, t = 1;
    for (i = 0; i < 8; i++)
        if (a[i] % 2 != 0) t *= 2;
    printf("%d %d\n", fun(a, 8), t);
    return 0;
}
\end{lstlisting}
\anslines{1.5cm}

\q （5 分）下列程序用于统计输入字符串中大写字母、小写字母、数字字符和其他字符的个数，请在标号 (1)、(2) 处填入正确内容，使程序完整正确。
\begin{lstlisting}[language={[custom]C}]
#include <stdio.h>
int main(void) {
    char s[80];
    int i, upper = 0, lower = 0, digit = 0, other = 0;
    gets(s);
    for (i = 0; (1) ; i++) {
        if (s[i] >= 'A' && s[i] <= 'Z') upper++;
        else if (s[i] >= 'a' && s[i] <= 'z') lower++;
        else if ( (2) ) digit++;
        else other++;
    }
    printf("upper=%d, lower=%d, digit=%d, other=%d\n",
           upper, lower, digit, other);
    return 0;
}
\end{lstlisting}
\anslines{1.0cm}
\end{qgroup}

\qtypename{四、程序设计题}{2 题，每题 5 分，共 10 分}
\begin{qgroup}
\q （5 分）从键盘输入 10 个整数存入一维数组，然后输出其中的最大值、最小值以及平均值（平均值保留两位小数）。要求写出完整的 C 程序。
\anslines{5.0cm}

\q （5 分）编写函数 \code{int is\_prime(int n)}：若正整数 \code{n} 是素数则返回 1，否则返回 0。再编写主函数，输出 100 到 200 之间的全部素数，要求每行输出 5 个。请写出完整的 C 程序。
\anslines{5.0cm}
\end{qgroup}

\end{document}
